% ==============================================================================
% 2. PHƯƠNG PHÁP VÀ NGUỒN THU THẬP DỮ LIỆU
% ==============================================================================
\needspace{0.5\textheight}
\section{Phương pháp và nguồn thu thập dữ liệu}

\subsection{Nguồn dữ liệu}

Bảng~\ref{tab:sources} lần lượt trả lời các câu hỏi về nguồn dữ liệu trong khung báo cáo.

\begin{table}[H]
\centering
\caption{Trả lời các câu hỏi về nguồn dữ liệu của khung báo cáo}
\label{tab:sources}
\begin{tabularx}{\textwidth}{>{\raggedright\arraybackslash}p{0.27\linewidth}L}
\toprule
\textbf{Câu hỏi} & \textbf{Trả lời} \\
\midrule
Dựa trên dataset có sẵn? & Không. Nhóm tự xây dựng bộ dữ liệu từ ảnh gốc. Bộ dữ liệu gần nhất là Vietnamese Sidewalk Violation VQA~\cite{dao2026viahevqa}, nhưng bộ này không có tọa độ và thời gian chụp nên không thể chia theo tuyến đường; nhóm không sử dụng ảnh hay nhãn của bộ này. \\
Từ API? & Có. Ảnh $360^\circ$ và metadata được lấy qua Mapillary Graph API~\cite{mapillary2024intro}; mạng lưới đường và các điểm ngữ cảnh (chợ, trường học) được lấy từ OpenStreetMap (OSM)~\cite{osm2026} qua Overpass API. \\
Web scraping? & Không. \\
Khảo sát, biểu mẫu? & Không. Nhãn do 5 thành viên nhóm gán trên Label Studio~\cite{labelstudio2020} (mục~\ref{sec:labeling}). \\
Giấy phép & Ảnh Mapillary theo giấy phép CC BY-SA 4.0, ghi công \cc. Phần làm mờ khuôn mặt và biển số của Mapillary được giữ nguyên. Ảnh không được đưa lên git. \\
Ảnh nhóm tự chụp & \textbf{Chưa có} (đề cương dự kiến ít nhất 100 ảnh). \\
\bottomrule
\end{tabularx}
\end{table}

Khu vực nghiên cứu là trung tâm TP.HCM (Quận 1, 3, 5, 10, Phú Nhuận và Bình Thạnh cũ), trong khung tọa độ 10,745--10,815$^\circ$ Bắc, 106,650--106,725$^\circ$ Đông. Trong khung này, Mapillary có 652.159 ảnh $360^\circ$ thuộc 1.234 chuỗi ảnh (quét ngày 22/10/2026), chủ yếu của ba tài khoản \lab{bemaps3\_sg}, \lab{phoenix\_mapping\_front} và \lab{phoenixmap\_front}, chụp trong các năm 2016, 2023 và 2024.

\subsection{Quy trình thu thập từng bước}
\label{sec:pipeline}

Bảng~\ref{tab:steps} liệt kê chín bước thu thập và kết quả của từng bước. Mỗi bước do một script riêng của nhóm thực hiện và được ghi lại trong nhật ký thu thập.

\begin{table}[H]
\centering
\caption{Các bước thu thập bộ chính và pilot}
\label{tab:steps}
\begin{tabularx}{\textwidth}{c >{\raggedright\arraybackslash}p{0.34\linewidth} L}
\toprule
\textbf{\#} & \textbf{Công việc} & \textbf{Kết quả} \\
\midrule
1 & Lấy các đường có tên và các điểm ngữ cảnh từ OSM & 7.283 đoạn đường, 6.715 điểm \\
2 & Quét metadata ảnh $360^\circ$ theo các ô lưới 0,001$^\circ$ & 652.159 ảnh; 2 ô nhỏ ở góc khung bị lỗi \\
3 & Gắn ảnh vào đường: cách đường $\le 20$ m và hướng di chuyển lệch trục đường $\le 30^\circ$; loại hẻm & 81,3\% số ảnh gắn được vào một đường \\
4 & Xem trước 56 đường, mỗi đường 20 vị trí, để chọn đường & Loại 7 đường ven kênh, cầu, đường hầm; chọn 17 đường \\
5 & Chia 17 đường thành train/val/test \emph{trước khi} gán nhãn & 10 đường train, 3 đường val, 4 đường test \\
6 & Pilot: 150 ảnh trên 3 đường train, cả 5 người cùng gán mỗi ảnh & Fleiss $\kappa$~\cite{fleiss1971kappa} 0,57; viết lại hướng dẫn gán nhãn (v1) \\
7 & Thu thập bộ chính: các vị trí cách nhau 25 m, tối đa 2 ảnh mỗi vị trí từ hai ngày khác nhau, tối đa 420 ảnh mỗi đường, bỏ các vị trí của pilot & Chọn 3.057 ảnh, loại 58 ảnh chụp ban đêm, còn 2.999 ảnh \\
8 & Gán nhãn bộ chính, đo đồng thuận, phân xử (mục~\ref{sec:labeling}) & Nhãn cuối cho 3.149 ảnh \\
9 & Bổ sung ảnh cho train từ 75 đường mới (mục~\ref{sec:enrich}) & 2.000 ảnh, trong đó 577 ảnh \lab{Bị chặn} \\
\bottomrule
\end{tabularx}
\end{table}

\textbf{Cắt ảnh.} Mỗi ảnh $360^\circ$ cho đúng một ảnh cắt: nhìn về bên phải xe, lệch $60^\circ$ so với hướng di chuyển (góc yaw), nghiêng xuống $4^\circ$ (góc pitch), trường nhìn $90^\circ$, kích thước $640\times640$ điểm ảnh, lưu dạng JPEG chất lượng 90. Hai quyết định dưới đây dựa trên phép đo thăm dò ngày 23/10/2026, không dựa trên nhãn:
\begin{itemize}
  \item \emph{Chỉ cắt ảnh bên phải.} Trên 1.166 ảnh $360^\circ$, một phép đo tự động (với tiêu chí chặt) phát hiện vỉa hè gần ở 38\% ảnh bên phải và 7\% ảnh bên trái. Ảnh bên trái phần lớn chỉ thấy lòng đường, dải phân cách hoặc vỉa hè phía đối diện. Hệ quả là vỉa hè bên trái của đường một chiều ít được quan sát.
  \item \emph{Góc lệch $60^\circ$ và góc nghiêng $4^\circ$} cho tỉ lệ quan sát được vỉa hè gần cao nhất trong 6 tổ hợp đã thử (góc lệch 60/75/90$^\circ$, góc nghiêng 4/10$^\circ$).
\end{itemize}

\textbf{Chia tập theo tuyến đường.} Mỗi đường thuộc trọn một tập, nhờ đó mô hình được đánh giá trên các đường chưa từng thấy. Bốn đường test nằm ở bốn khu vực khác nhau. Danh sách đường và tập được xác định ngày 23/10/2026, trước khi có bất kỳ nhãn nào của bộ chính. Hiện việc chia tập đã hoàn tất nhưng \emph{chưa được khóa}, do nhóm còn chờ quyết định cách xử lý các ảnh train nằm gần giao lộ với val/test. Nhóm giữ nguyên các đường test đã định: việc test chỉ có 47 mẫu dương là một hạn chế về độ chính xác của đánh giá, không phải tiêu chí để thay đổi tập dựa trên nhãn đã quan sát.

\subsection{Gán nhãn}
\label{sec:labeling}

Nhãn do người gán, theo hướng dẫn gán nhãn phiên bản v1 (bản đã được nhóm thống nhất nguyên văn). Mỗi ảnh được xét qua ba câu hỏi theo thứ tự:
\begin{enumerate}
  \item Ảnh có quan sát được không? Ảnh mờ, tối hoặc cháy sáng đến mức không thể đánh giá thì gán \lab{I}.
  \item Có thấy vỉa hè gần không? Phải có một đoạn mép vỉa hè nằm dưới vạch 3/5 chiều cao ảnh; nếu không (chỉ thấy lòng đường, vỉa hè ở xa, giao lộ, đầu hẻm, hoặc bị che quá nửa) thì gán \lab{I}.
  \item Lối đi có liên tục không? Nếu không thể đi thẳng qua thì gán \lab{1}. Người và xe đang di chuyển không được tính là vật cản; chỉ xét vật cản trong phạm vi tới khoảng giữa ảnh.
\end{enumerate}

Ngoài nhãn chính, người gán ghi lý do khi chọn \lab{I} (4 lý do), có thể đánh dấu loại vật cản (6 loại, không bắt buộc), ô \lab{Khó / không chắc} và ghi chú tự do. Giao diện gán nhãn trên Label Studio được cấu hình theo đúng ba câu hỏi trên.

\begin{table}[H]
\centering
\caption{Phân công gán nhãn}
\label{tab:assignment}
\begin{tabularx}{\textwidth}{>{\raggedright\arraybackslash}p{0.22\linewidth} r L}
\toprule
\textbf{Phần} & \textbf{Ảnh} & \textbf{Cách gán} \\
\midrule
Test & 674 & Hai người (Học, Loan) gán độc lập mọi ảnh \\
Val & 486 & Hai người (Minh, Khải) gán độc lập mọi ảnh \\
Train, bộ chính & 1.839 & Mỗi ảnh một người; 276 ảnh (15\%) được người thứ hai gán chéo \\
Pilot & 150 & 5 người gán theo hướng dẫn v0; sau đó cả 150 ảnh được xem lại theo v1 \\
Ảnh bổ sung & 2.000 & Mỗi ảnh một người (5 người, mỗi người 400 ảnh được chia ngẫu nhiên); 203 ảnh được người thứ hai kiểm tra chéo \\
\bottomrule
\end{tabularx}
\end{table}

Bảng~\ref{tab:assignment} cho biết phân công gán nhãn. Người gán không thấy tên đường, tập dữ liệu, điểm của mô hình hay phiếu của người khác. Ảnh của các tập được trộn lẫn trong cùng một dự án gán nhãn. Bộ chính có tổng cộng $5\times887 = 4.435$ lượt gán, không thiếu lượt nào. Kết quả đồng thuận và phân xử được trình bày ở mục~\ref{sec:qa}.

\subsection{Bổ sung mẫu dương cho train}
\label{sec:enrich}

Sau đợt gán nhãn đầu tiên, train chỉ có 120 mẫu dương (gồm 14 ảnh pilot). Nhóm bổ sung ảnh \emph{chỉ vào train}, theo hướng chủ động tìm mẫu dương của lớp hiếm~\cite{attenberg2010search}:

\begin{enumerate}
  \item \textbf{Bộ chọn ảnh.} Hồi quy logistic trên vector đặc trưng (embedding) DINOv2 ViT-B/14~\cite{oquab2023dinov2}, huấn luyện trên 1.989 nhãn train (120 mẫu dương). Bộ chọn này được lựa chọn qua một phép so sánh chỉ dùng nhãn train, theo cách để riêng từng tuyến đường (leave-one-street-out): trong 20\% ảnh có điểm cao nhất có 22,0\% ảnh \lab{Bị chặn}, so với 13,9\% của bộ chấm điểm cũ (chênh lệch $+8{,}2$ điểm phần trăm, khoảng tin cậy 95\% $[4{,}3;\,11{,}7]$) và 5,8\% nếu lấy mẫu ngẫu nhiên.
  \item \textbf{Chọn đường.} Quét 150 đường chưa sử dụng (mỗi đường 30 ảnh), xếp hạng theo điểm trung bình và lấy 75 đường đứng đầu. Đường trùng tên với đường val/test bị loại, kể cả các đoạn khác của cùng đường.
  \item \textbf{Chọn ảnh.} Kho ứng viên gồm 5.035 ảnh. Nhóm chọn 2.000 ảnh: 1.400 ảnh có điểm cao nhất (\emph{tầng điểm cao}, điểm $\ge 0{,}61$) và 600 ảnh lấy ngẫu nhiên trong phần còn lại (\emph{tầng ngẫu nhiên}). Các vị trí cách ảnh val/test dưới 40 m bị loại.
  \item \textbf{Gán nhãn.} Giống phần train của bộ chính, mỗi ảnh do một người gán; người gán không biết ảnh thuộc tầng nào.
\end{enumerate}

Điểm của bộ chọn ảnh \emph{chỉ dùng để chọn ảnh cho người gán, không bao giờ được dùng làm nhãn}. Vì ảnh bổ sung được chọn có chủ đích, phân phối của chúng lệch khỏi phân phối tự nhiên~\cite{farquhar2021bias}. Tầng ngẫu nhiên (30\%) được lấy từ phần còn lại sau khi chọn tầng điểm cao, nên không phải là mẫu ngẫu nhiên vô điều kiện của toàn kho ứng viên; tầng này được giữ để so sánh hiệu suất thu được mẫu dương trên đường mới. Mọi thống kê trong báo cáo đều tách ảnh bổ sung khỏi bộ chính.

\subsection{Công cụ và định dạng lưu trữ}

Bảng~\ref{tab:tools} tóm tắt công cụ và định dạng lưu trữ ở từng bước; vị trí lưu từng loại dữ liệu được trình bày ở mục~\ref{sec:storage}.

\begin{table}[H]
\centering
\caption{Công cụ sử dụng và định dạng lưu trữ}
\label{tab:tools}
\begin{tabularx}{\textwidth}{>{\raggedright\arraybackslash}p{0.25\linewidth} >{\raggedright\arraybackslash}p{0.30\linewidth} L}
\toprule
\textbf{Công việc} & \textbf{Công cụ} & \textbf{Định dạng lưu} \\
\midrule
Mạng lưới đường, điểm ngữ cảnh & Overpass API, shapely & JSON \\
Metadata ảnh & Mapillary Graph API, requests, pandas & Parquet \\
Ảnh $360^\circ$ gốc & Mapillary (ảnh thu nhỏ 2.048 điểm ảnh) & JPEG \\
Cắt ảnh & py360convert, Pillow & JPEG $640\times640$ \\
Danh mục ảnh (manifest) & pandas & CSV UTF-8, mỗi ảnh một dòng \\
Gán nhãn & Label Studio tự triển khai trên AWS (Docker, Caddy) & Export JSON trên server; bảng nhãn CSV trong repo \\
Đo đồng thuận, thống kê & pandas, SciPy, scikit-learn & CSV, Markdown \\
Bộ chọn ảnh, đo độ phủ vỉa hè (chỉ phục vụ chọn mẫu) & PyTorch, Transformers (DINOv2~\cite{oquab2023dinov2}, CLIP~\cite{radford2021clip}, SegFormer~\cite{xie2021segformer}) & Embedding NPZ \\
Cấu hình đường và tập & & YAML \\
EDA của báo cáo này & pandas, SciPy, Matplotlib & CSV, JSON, PDF, JPEG \\
\bottomrule
\end{tabularx}
\end{table}
